% Part: counterfactuals
% Chapter: introduction
% Section: paradoxes-material

\documentclass[../../../include/open-logic-section]{subfiles}

\begin{document}

\olfileid{cnt}{int}{par}

\olsection{في مفارقات الشرط المادي}

كان سي.~آي. لويس في طليعة من أنكروا أن يكون~$!A \lif !B$ تمثيلًا
صحيحًا لكل قول إنجليزي من نحو «إذا \dots، فإن \dots». وخصّ بالنقد
استعمال الشرط المادي في كتاب وايتهد وراسل \emph{مبادئ الرياضيات}، إذ
كانا يقرآن~$\lif$ «يستلزم». واعتراضه في موضعه إن حُمّل~$\lif$ هذا
المعنى: فكل قضية كاذبة~$\OLId{latin-ordinary}{p}$ تستلزم عندئذ أن~$\OLId{latin-ordinary}{p}$ تستلزم~$\OLId{latin-ordinary}{q}$، لأن
$\OLId{latin-ordinary}{p} \lif (\OLId{latin-ordinary}{p} \lif \OLId{latin-ordinary}{q})$ تصدق إذا كذبت~$\OLId{latin-ordinary}{p}$؛ وكل قضية صادقة~$\OLId{latin-ordinary}{q}$ تستلزم
كذلك أن~$\OLId{latin-ordinary}{p}$ تستلزم~$\OLId{latin-ordinary}{q}$، لأن~$\OLId{latin-ordinary}{q} \lif (\OLId{latin-ordinary}{p} \lif \OLId{latin-ordinary}{q})$ تصدق إذا صدقت~$\OLId{latin-ordinary}{q}$.

والمعلوم عند المناطقة أن \emph{الاستلزام} المنطقي ليس رابطًا، بل علاقة
بين !!{formula} أو العبارات. فلا ينبغي إذن أن نقرأ
$\lif$ بمعنى «يستلزم»، اتقاءً للالتباس.\footnote{ما زالت قراءة
  «$\lif$» بمعنى «يستلزم» واسعة الانتشار بين الرياضيين وعلماء الحاسوب،
  مع أن الفلاسفة يتقون اللبس الذي كشفه لويس بقراءته «فقط إذا».} فإذا
تحاشينا تلك القراءة زال أصل اعتراض لويس: فليست~$\OLId{latin-ordinary}{p}$ «مستلزمة» لـ~$\OLId{latin-ordinary}{q}$
لمجرد أن الرمز~$\OLId{latin-ordinary}{p}$ ناب عن جملة إنجليزية كاذبة في الواقع. والحكم في
$\OLId{latin-ordinary}{p} \Entails \OLId{latin-ordinary}{q}$ مردّه إلى \emph{جميع} !!{valuation}؛ ولذلك يبقى
$\OLId{latin-ordinary}{p} \Entails/ \OLId{latin-ordinary}{q}$ وإن صادف أن الجملة الممثلة بـ~$\OLId{latin-ordinary}{p}$ كاذبة.

على أن الغرابة تبقى في بعض عبارات «إذا \dots، فإن \dots»، ومنها مثال
لويس:
\begin{quote}
إذا كان القمر مصنوعًا من الجبن الأخضر، فإن~$\OLMathNumeral{2}+\OLMathNumeral{2}=\OLMathNumeral{4}$.
\end{quote}
وكذلك في الاستدلالين:
\begin{quote}
  القمر ليس مصنوعًا من الجبن الأخضر. ولذلك، إذا كان القمر مصنوعًا
  من الجبن الأخضر، فإن~$\OLMathNumeral{2}+\OLMathNumeral{2}=\OLMathNumeral{4}$.

  $\OLMathNumeral{2}+\OLMathNumeral{2} = \OLMathNumeral{4}$. ولذلك، إذا كان القمر مصنوعًا
  من الجبن الأخضر، فإن~$\OLMathNumeral{2}+\OLMathNumeral{2}=\OLMathNumeral{4}$.
\end{quote}
ومع ذلك، إن لم تكن «إذا \dots، فإن \dots» إلا~$\lif$، صدقت الجملة
ولزم صحة الاستدلالين من غير قيد.

ومثال آخر ناشئ من القضية الصادقة منطقيًا
$(!A \lif !B) \lor (!B \lif !A)$. فإن أخذنا اعتباطًا جملتين خبريتين
$\OLId{latin-looped}{S}$ و$\OLId{latin-looped}{T}$ من صحيفة، قضت هذه القضية بصدق القول: «إذا كانت~$\OLId{latin-looped}{S}$ فإن~$\OLId{latin-looped}{T}$،
أو إذا كانت~$\OLId{latin-looped}{T}$ فإن~$\OLId{latin-looped}{S}$».

\end{document}



